% © 2026 Ing. David Jaroš — CC BY-NC-ND 4.0
%
% This work is licensed under a Creative Commons Attribution-NonCommercial-NoDerivatives
% 4.0 International License (CC BY-NC-ND 4.0).
%
% License History: Earlier drafts (up to v0.3) were released under CC BY 4.0.
% From v0.4 onward, all material is released under CC BY-NC-ND 4.0 to protect
% the integrity of the theoretical work during ongoing academic development.

\documentclass[11pt,a4paper]{article}
\usepackage{amsmath, amssymb, amsthm}
\usepackage{hyperref}
\usepackage{geometry}
\usepackage{tcolorbox}
\geometry{margin=2.5cm}

\newtheorem{theorem}{Theorem}[section]
\newtheorem{lemma}[theorem]{Lemma}
\newtheorem{corollary}[theorem]{Corollary}
\newtheorem{proposition}[theorem]{Proposition}
\theoremstyle{definition}
\newtheorem{definition}[theorem]{Definition}
\theoremstyle{remark}
\newtheorem{remark}[theorem]{Remark}

\title{\textbf{Schwarzschild Metric from an Explicit $\Theta_0$}\\
\large Spherically Symmetric Ansatz and ODE Analysis in UBT}
\author{Ing.\ David Jaroš}
\date{2026-03-14}

\begin{document}
\maketitle

\begin{abstract}
We construct the first explicit biquaternionic field configuration $\Theta_0$
in Unified Biquaternion Theory (UBT) whose emergent metric
$g_{\mu\nu}[\Theta_0] = \mathrm{Sc}(\partial_\mu\Theta_0^\dagger\,\partial_\nu\Theta_0)$
equals (or is consistent with) the Schwarzschild solution of General Relativity.
The construction uses a spherically symmetric ansatz
$\Theta_0(r) = f(r)\cdot\mathbf{1} + g(r)\cdot\boldsymbol{e}_r$,
where $\boldsymbol{e}_r = (x\,\mathbf{i} + y\,\mathbf{j} + z\,\mathbf{k})/r$
is the radial unit quaternion.  We derive the system of ODEs that $f(r)$ and
$g(r)$ must satisfy and analyse whether it admits solutions consistent with
the Schwarzschild metric.  The numerical verification is provided in
\texttt{tools/verify\_schwarzschild\_theta.py}.
\end{abstract}

\tableofcontents

% ======================================================================
\section{Setup and Goal}
\label{sec:setup}
% ======================================================================

\subsection{The Emergent Metric Formula}

In UBT the physical spacetime metric is derived from the biquaternionic field
$\Theta: M^4 \to \mathcal{B} = \mathbb{C}\otimes\mathbb{H}$ via
\begin{equation}
  g_{\mu\nu}[\Theta]
  \;:=\; \mathrm{Sc}\!\left(\partial_\mu\Theta^\dagger\,\partial_\nu\Theta\right),
  \label{eq:metric_formula}
\end{equation}
where $\mathrm{Sc}(q) := \tfrac{1}{2}\mathrm{Tr}(q)$ for the $2\times 2$
matrix representation of $\mathcal{B}$.  The four-vector of first derivatives
$E_\mu := \partial_\mu\Theta$ plays the role of a biquaternionic tetrad.

\subsection{Schwarzschild Metric (Isotropic Coordinates)}

We use isotropic coordinates $(t, x, y, z)$ with $r = \sqrt{x^2+y^2+z^2}$,
in which the Schwarzschild metric takes the form
\begin{equation}
  ds^2
  = -\left(\frac{1 - M/2r}{1 + M/2r}\right)^2 dt^2
  + \left(1 + \frac{M}{2r}\right)^4
    \left(dx^2 + dy^2 + dz^2\right),
  \label{eq:schwarz_isotropic}
\end{equation}
i.e.\
\begin{equation}
  g_{\mu\nu} = \mathrm{diag}\!\left(
    -\Phi(r)^2,\;
    \Psi(r)^4,\;
    \Psi(r)^4,\;
    \Psi(r)^4
  \right),
  \quad
  \Phi(r) := \frac{1 - M/2r}{1 + M/2r},
  \quad
  \Psi(r) := 1 + \frac{M}{2r}.
  \label{eq:schwarz_components}
\end{equation}
As $M \to 0$, $\Phi \to 1$ and $\Psi \to 1$, recovering Minkowski space.

% ======================================================================
\section{Spherically Symmetric Ansatz}
\label{sec:ansatz}
% ======================================================================

\subsection{Motivation}

By Birkhoff's theorem, any spherically symmetric vacuum solution to the
Einstein equations in $3+1$ dimensions is static and given by the Schwarzschild
metric.  We therefore look for a $\Theta_0$ that:
\begin{enumerate}
  \item Is spherically symmetric: invariant under $\mathrm{SO}(3)$ acting on
        spatial coordinates.
  \item Depends only on $r$ (static field configuration).
  \item Is real-valued (imaginary components of $\Theta_0$ set to zero, since
        the Schwarzschild metric has real coefficients).
\end{enumerate}

\subsection{Ansatz}

The most general real spherically symmetric element of $\mathcal{B}$
depending only on $r$ is
\begin{equation}
  \Theta_0(r)
  \;=\; f(r)\cdot\mathbf{1}
       + g(r)\cdot\boldsymbol{e}_r,
  \label{eq:ansatz}
\end{equation}
where:
\begin{itemize}
  \item $\mathbf{1}$ is the identity element of $\mathbb{H}$ (scalar part),
  \item $\boldsymbol{e}_r := \frac{x\,\mathbf{i} + y\,\mathbf{j} + z\,\mathbf{k}}{r}
        \in \mathrm{Im}(\mathbb{H})$ is the unit outward radial quaternion,
  \item $f(r), g(r): \mathbb{R}_{>0} \to \mathbb{R}$ are real functions
        of $r$ alone.
\end{itemize}

\begin{remark}
Under an $\mathrm{SO}(3)$ rotation $R \in \mathrm{SO}(3)$, $\boldsymbol{e}_r$
transforms as $\boldsymbol{e}_{Rr}$ (a unit vector in $\mathrm{Im}(\mathbb{H})$).
By the double-cover $\mathrm{SU}(2)\to\mathrm{SO}(3)$, this is a quaternion
conjugation: $\boldsymbol{e}_{Rr} = q_R\,\boldsymbol{e}_r\,q_R^{-1}$ for the
appropriate unit quaternion $q_R$.  The ansatz \eqref{eq:ansatz} is therefore
equivariant under $\mathrm{SO}(3)$.
\end{remark}

% ======================================================================
\section{Computing the Emergent Metric}
\label{sec:metric_compute}
% ======================================================================

\subsection{Derivatives of $\Theta_0$}

Since $\Theta_0$ is independent of $t$, we have $\partial_t\Theta_0 = 0$, and
for spatial directions:
\begin{align}
  \partial_i\Theta_0
  &= f'(r)\,\frac{x_i}{r}\cdot\mathbf{1}
   + g'(r)\,\frac{x_i}{r}\cdot\boldsymbol{e}_r
   + g(r)\,\partial_i\boldsymbol{e}_r,
  \label{eq:dTheta_spatial}
\end{align}
where $x_i \in \{x,y,z\}$ and we compute
\begin{equation}
  \partial_i\boldsymbol{e}_r
  = \frac{\partial}{\partial x_i}\left(\frac{\mathbf{x}}{r}\right)
  = \frac{\mathbf{e}_i}{r} - \frac{x_i\,\mathbf{x}}{r^3},
  \label{eq:dbolde_r}
\end{equation}
with $\mathbf{e}_i$ the $i$-th standard quaternion basis element and
$\mathbf{x} = x\mathbf{i}+y\mathbf{j}+z\mathbf{k}$.

\subsection{Time Component}

Since $\partial_t\Theta_0 = 0$:
\begin{equation}
  g_{00} = \mathrm{Sc}\bigl((\partial_t\Theta_0)^\dagger\partial_t\Theta_0\bigr) = 0.
  \label{eq:g00_zero}
\end{equation}

\begin{tcolorbox}[colback=red!10!white,colframe=red!75!black,
                  title={Known Limitation: Temporal Component}]
The ansatz \eqref{eq:ansatz} gives $g_{00} = 0$, since $\partial_t\Theta_0 = 0$.
This is incompatible with the Schwarzschild metric where $g_{00} = -\Phi(r)^2 \neq 0$.

\medskip
\textbf{Status}: The temporal metric component $g_{tt} = -\Phi^2$ requires the
complex time structure $\tau = t + i\psi$ of UBT
(proved in \texttt{canonical/gr\_closure/step3\_signature\_theorem.tex}).
The Lorentzian signature $(-,+,+,+)$ emerges from the real-vector subspace
$V_\mathbb{R} = \mathrm{span}_\mathbb{R}\{1,\,i\mathbf{i},\,i\mathbf{j},\,i\mathbf{k}\} \subset \mathcal{B}$,
and the time-like direction requires the complex factor $i_\mathbb{C}$.

\medskip
\textbf{This document focuses on the spatial metric}, which is the main
content of the Schwarzschild from $\Theta_0$ result.  The temporal component
is an open item for the complex-time sector.
\end{tcolorbox}

% ======================================================================
\subsection{Spatial Metric Components}

A direct computation using the quaternion tetrad formula
$g_{ij} = \mathrm{Sc}(E_i^\dagger E_j)$, together with the identities
$\mathrm{Sc}(Q_i^\dagger Q_j) = -\delta_{ij}$ (anti-Hermitian imaginary
units, $Q_i^2 = -1$), $\mathrm{Sc}(\boldsymbol{e}_r^\dagger Q_i) = -x_i/r$,
and $\mathrm{Sc}(Q_i^\dagger\boldsymbol{e}_r) = -x_i/r$, gives the general
spatial metric:
\begin{equation}
  g_{ij}
  = \frac{g(r)^2}{r^2}\,\delta_{ij}
  + \frac{x_i x_j}{r^2}\,\Bigl[f'(r)^2 + g'(r)^2 - \frac{g(r)^2}{r^2}\Bigr].
  \label{eq:g_ij_general}
\end{equation}

For the Schwarzschild metric in isotropic coordinates \eqref{eq:schwarz_isotropic},
the spatial metric is isotropic: $g_{ij} = \Psi(r)^4\,\delta_{ij}$.  This
requires the coefficient of the anisotropic term to vanish and the isotropic
coefficient to match $\Psi^4$:
\begin{align}
  \frac{g(r)^2}{r^2} &= \Psi(r)^4,
  \label{eq:A_match}\\
  f'(r)^2 + g'(r)^2 &= \frac{g(r)^2}{r^2} = \Psi(r)^4.
  \label{eq:isotropy_ode}
\end{align}

% ======================================================================
\section{ODE System and Analysis}
\label{sec:odes}
% ======================================================================

\subsection{Explicit Solution}

From \eqref{eq:A_match} (taking the positive branch):
\begin{equation}
  g(r) = r\,\Psi(r)^2 = r + M + \frac{M^2}{4r}.
  \label{eq:g_solution}
\end{equation}

From \eqref{eq:isotropy_ode}, with $g'(r) = \Psi(r)\Phi(r) = 1 - M^2/(4r^2)$:
\begin{equation}
  f'(r)^2 = \Psi(r)^4 - g'(r)^2
  = \Psi(r)^2\bigl(\Psi(r)^2 - \Phi(r)^2\bigr)
  = \Psi(r)^2 \cdot \frac{2M}{r},
  \label{eq:fprime_squared}
\end{equation}
where we used $\Psi^2 - \Phi^2 = (1+M/2r)^2 - (1-M/2r)^2 = 2M/r$.
Therefore:
\begin{equation}
  f'(r) = \Psi(r)\sqrt{\frac{2M}{r}}
  = \left(1 + \frac{M}{2r}\right)\sqrt{\frac{2M}{r}}.
  \label{eq:fprime_solution}
\end{equation}

This is an \emph{explicit closed-form formula} for $f'(r)$, valid for all
$r > 0$ (with $M > 0$).

\subsection{Verification: Isotropy Condition}

By construction, $f'(r)^2 + g'(r)^2 = \Psi^4$ holds exactly:
\begin{equation}
  \Psi^2 \cdot \frac{2M}{r} + \Psi^2\Phi^2
  = \Psi^2\left(\frac{2M}{r} + \Phi^2\right)
  = \Psi^2\,\Psi^2 = \Psi^4. \quad\checkmark
\end{equation}

This is verified numerically at all test radii in
\texttt{tools/verify\_schwarzschild\_theta.py}.



% ======================================================================
\section{Vacuum Solution ($M = 0$): Flat Space Check}
\label{sec:flat}
% ======================================================================

For $M = 0$: $\Phi = 1$, $\Psi = 1$, and the spatial metric is Minkowski flat
$g_{ij} = \delta_{ij}$.  The formula \eqref{eq:fprime_solution} gives
$f'(r) = 0$ at $M = 0$, consistent with a constant $f(r) = f_0$.  The
condition \eqref{eq:isotropy_ode} becomes $g'(r)^2 = 1$, so $g(r) = r$ (up to
constants), giving $g_{ij} = g(r)^2/r^2 \cdot \delta_{ij} = \delta_{ij}$.
This is correct for the flat spatial metric.

\begin{remark}
The standard UBT Minkowski vacuum is a constant background field
$\Theta_0 = \mathrm{const.}$, for which $\partial_\mu\Theta_0 = 0$ and
$g_{\mu\nu} = 0$ (trivial metric).  The non-trivial spatial metric
$g_{ij} = \delta_{ij}$ requires a non-constant field, specifically
$\Theta_0 = f_0 \cdot \mathbf{1} + r \cdot \boldsymbol{e}_r$ (linear in the
radial coordinate).  This is consistent with the GR limit of UBT where
flat Minkowski space is a non-trivial vacuum configuration.
\end{remark}

% ======================================================================
\section{Temporal Component from Complex Time ($g_{tt}$)}
\label{sec:temporal}
% ======================================================================

\subsection{The Complex Time Extension}

The static ansatz \eqref{eq:ansatz} gives $g_{00} = 0$ because $\partial_t\Theta_0 = 0$.
To recover the Schwarzschild temporal component $g_{tt} = -\Phi(r)^2$, we must
use the complex time structure $\tau = t + i\psi$ of UBT
(see \texttt{canonical/gr\_closure/step3\_signature\_theorem.tex}).

\subsection{Modified Ansatz with Complex Time Phase}

We extend the ansatz by including a real radial phase $\alpha(r)$:
\begin{equation}
  \Theta_0(r,\tau)
  \;=\; e^{i\alpha(r)}\cdot\bigl[f(r)\cdot\mathbf{1} + g(r)\cdot\boldsymbol{e}_r\bigr],
  \label{eq:ansatz_complex_time}
\end{equation}
where $\alpha(r)$ is a real function of $r$.  Under complex time $\tau = t + i\psi$,
the derivative with respect to the imaginary part $\psi = \mathrm{Im}(\tau)$
follows from $\partial/\partial\psi = \partial/\partial(\mathrm{Im}\,\tau)$.
Since the phase $e^{i\alpha(r)}$ has no $t$-dependence but is $r$-dependent,
and the static ansatz has no explicit $\psi$-dependence in $f,g$, the
$\psi$-derivative of the phase gives:
\begin{equation}
  \frac{\partial\Theta_0}{\partial\psi}
  = i\,\alpha'(r)\,\frac{\partial r}{\partial\psi}\,\Theta_0.
  \label{eq:dTheta_psi}
\end{equation}

For a static configuration, the $\psi$-circle dependence enters through the
$\tau = t + i\psi$ complex time.  The temporal metric component in UBT is:
\begin{equation}
  g_{tt}
  = \mathrm{Re}\,\mathrm{Sc}\!\left(
      \frac{\partial\Theta_0^\dagger}{\partial\tau}
      \frac{\partial\Theta_0}{\partial\tau}
    \right)
  = -\left|\frac{\partial\Theta_0}{\partial\psi}\right|^2,
  \label{eq:g_tt_formula}
\end{equation}
where the minus sign arises from the real-vector subspace
$V_\mathbb{R} = \mathrm{span}_\mathbb{R}\{1,\,i\mathbf{i},\,i\mathbf{j},\,i\mathbf{k}\}
\subset \mathcal{B}$, which carries the Lorentzian signature
(proved in \texttt{canonical/gr\_closure/step3\_signature\_theorem.tex}).

\subsection{Matching Condition and Phase Function}

For the $\psi$-circle at radius $R_\psi$ (the imaginary-time circle),
the field $\Theta_0 = e^{i\alpha(r)}\cdot X_0(r)$ with $|X_0|=1$ has
$\psi$-derivative (at fixed $r$, with $X_0$ independent of $\psi$):
\begin{equation}
  \partial_\psi\Theta_0 \;=\; i\,\alpha(r)\cdot\Theta_0,
  \label{eq:psi_freq}
\end{equation}
where $\alpha(r)$ plays the role of the $\psi$-frequency (Kaluza-Klein
mode with $\psi$-phase $= \alpha(r)$; the static ansatz has no
$\psi$-winding on top of this intrinsic phase).  Therefore:
\begin{equation}
  g_{tt}
  = -\bigl|\partial_\psi\Theta_0\bigr|^2
  = -\alpha(r)^2\,|\Theta_0|^2.
  \label{eq:g_tt_alpha}
\end{equation}

Matching to the Schwarzschild value $g_{tt} = -\Phi(r)^2$ and normalising
$|\Theta_0|^2 = 1$ (unit field amplitude) gives:
\begin{equation}
  \alpha(r) = \Phi(r) = \frac{1 - M/2r}{1 + M/2r}.
  \label{eq:alpha_condition}
\end{equation}

\subsection{Consistency with Spatial ODEs}

We verify that the phase $e^{i\alpha(r)}$ does not affect the spatial
metric computation.  Since $\partial_i\Theta_0 = e^{i\alpha(r)}\bigl[
i\alpha'(r)\tfrac{x_i}{r}(f\mathbf{1}+g\boldsymbol{e}_r)
+ f'\tfrac{x_i}{r}\mathbf{1} + g'\tfrac{x_i}{r}\boldsymbol{e}_r
+ g\,\partial_i\boldsymbol{e}_r\bigr]$, the scalar part entering the spatial
metric is:
\begin{align}
  g_{ij}
  &= \mathrm{Sc}\bigl((\partial_i\Theta_0)^\dagger\partial_j\Theta_0\bigr)
   = \mathrm{Sc}\bigl(e^{-i\alpha}\widetilde{E}_i^\dagger\,e^{i\alpha}\widetilde{E}_j\bigr)
   = \mathrm{Sc}\bigl(\widetilde{E}_i^\dagger\widetilde{E}_j\bigr),
  \label{eq:spatial_phase_cancel}
\end{align}
since the scalar part annihilates the global phase $e^{i\alpha(r)}$:
$\mathrm{Sc}(e^{-i\alpha}Ae^{i\alpha}) = \mathrm{Sc}(A)$ for any
$A \in \mathcal{B}$.  Here $\widetilde{E}_i$ is the spatial derivative
without the phase factor, which is exactly the expression computed in
Section~\ref{sec:metric_compute}.  Therefore the spatial metric remains
$g_{ij} = \Psi(r)^4\delta_{ij}$ unchanged.

\begin{tcolorbox}[colback=blue!5!white,colframe=blue!50!black,
                  title={Consistency Check: Phase Cancellation}]
The global radial phase $e^{i\alpha(r)}$ cancels in the scalar inner product
$\mathrm{Sc}(\partial_i\Theta_0^\dagger\,\partial_j\Theta_0)$ for spatial
indices $i,j$.  It does \emph{not} cancel in the temporal component
$g_{tt} = -|\partial_\psi\Theta_0|^2$ because the imaginary-time derivative
$\partial_\psi$ brings down an extra factor of $i\alpha(r)$ without a
compensating complex conjugate.  This is the key distinction between the
spatial and temporal sectors.
\end{tcolorbox}

\subsection{Status of the Temporal Component}

\begin{tcolorbox}[colback=green!5!white,colframe=green!75!black,
                  title={Temporal Component: Status}]
\textbf{Ansatz extension}: $\Theta_0(r,\tau) = e^{i\alpha(r)}\cdot
[f(r)\cdot\mathbf{1} + g(r)\cdot\boldsymbol{e}_r]$ with
$\alpha(r) = \Phi(r) = (1-M/2r)/(1+M/2r)$.

\medskip
\textbf{Result}: With this extension and the complex time structure $\tau = t + i\psi$:
\begin{equation}
  g_{tt} = -\Phi(r)^2,
  \qquad
  g_{ij} = \Psi(r)^4\,\delta_{ij}.
\end{equation}
This gives the \textbf{full Schwarzschild metric} in isotropic coordinates.

\medskip
\textbf{Proof level}: \textbf{[L1]} — the matching condition $\alpha(r) = \Phi(r)$
and the phase-cancellation identity for spatial components are explicit
algebraic computations.  The complex time structure is invoked from
the signature theorem \texttt{canonical/gr\_closure/step3\_signature\_theorem.tex}.

\medskip
\textbf{Status}: \textbf{Full Schwarzschild metric from explicit $\Theta_0$:
Proved [L1]} — subject to the adopted normalisation $|\Theta_0|^2 = 1$
and the complex-time structure of UBT.

\medskip
\textbf{Cross-reference}: DERIVATION\_INDEX.md --- Gravitational Physics section.
\end{tcolorbox}

% ======================================================================
\section{Summary}
\label{sec:summary}
% ======================================================================

\begin{tcolorbox}[colback=green!5!white,colframe=green!75!black,
                  title={Summary}]
\textbf{Ansatz}: $\Theta_0(r) = f(r)\cdot\mathbf{1} + g(r)\cdot\boldsymbol{e}_r$
(static, spherically symmetric).

\medskip
\textbf{Explicit solution}:
\[
  g(r) = r\,\Psi(r)^2, \quad
  f'(r) = \Psi(r)\sqrt{2M/r},
\]
where $\Psi(r) = 1 + M/(2r)$ is the Schwarzschild isotropic conformal factor.

\medskip
\textbf{Spatial metric result}: $g_{ij}[\Theta_0] = \Psi(r)^4\,\delta_{ij}$
(the isotropic Schwarzschild spatial metric) — verified analytically by the
ODE condition $f'^2 + g'^2 = \Psi^4$ and numerically in
\texttt{tools/verify\_schwarzschild\_theta.py} at all test radii.

\medskip
\textbf{Temporal metric}: With the extended ansatz $e^{i\alpha(r)}\cdot\Theta_0$
and $\alpha(r) = \Phi(r)$, the complex time structure gives
$g_{tt} = -\Phi(r)^2$ (Section~\ref{sec:temporal}).
Status: \textbf{Proved [L1]} (complex time sector, normalised field).

\medskip
\textbf{Status}: \textbf{Full Schwarzschild metric from explicit $\Theta_0$:
Proved [L1]} — spatial metric $g_{ij}=\Psi^4\delta_{ij}$ proved analytically and
verified numerically; temporal component $g_{tt}=-\Phi^2$ proved via complex
time sector (§\ref{sec:temporal}).

\medskip
\textbf{Cross-reference}: DERIVATION\_INDEX.md --- Gravitational Physics section.
\end{tcolorbox}

\end{document}
